\section{Related Work}

The landscape of cache replacement policies has evolved from simple
static heuristics to complex, adaptive algorithms, and recently, to
machine learning-augmented systems. Research typically optimizes for
two often conflicting goals: maximizing the hit ratio (scan
resistance) and minimizing implementation overhead (concurrency and
scalability).

\subsection{Classical Replacement and Scalability}
The Least Recently Used (LRU) algorithm has historically been the
baseline for replacement policies due to its exploitation of temporal
locality \cite{Denning2005Locality}. However, exact LRU requires a
doubly-linked list update on every cache access. In modern multicore
processors, the fine-grained locking required to protect this
structure results in severe contention, limiting system throughput
\cite{Herlihy2006The}.

To address scalability, the CLOCK algorithm approximates LRU using a
circular buffer and reference bits, eliminating the need for list
manipulation on hits. Eads et al. further refined this for
high-concurrency environments with \textbf{NbQ-CLOCK} (Non-blocking
Queue CLOCK) \cite{Eads2014NbQCLOCK}. NbQ-CLOCK replaces the circular
buffer with a concurrent, lock-free queue, allowing for high
throughput. However, both standard CLOCK and NbQ-CLOCK suffer from
``blind promotion,'' where any accessed item is granted a second
chance. This makes them susceptible to cache pollution during scan
operations (sequences of ``one-hit wonders''), which displace
high-utility data \cite{Jiang2005CLOCKPro}.

\subsection{Scan Resistance and Advanced Heuristics}
To mitigate cache pollution caused by scans and Zipfian long-tail
distributions \cite{Breslau1999Web}, several heuristics were
developed to separate frequency from recency.

\textbf{LIRS} (Low Inter-reference Recency Set) utilizes the
Inter-Reference Recency (IRR) between the last two accesses to
predict future utility, effectively filtering out items with weak
locality \cite{Jiang2002LIRS}. \textbf{ARC} (Adaptive Replacement
Cache) dynamically balances recency and frequency by maintaining two
lists and adapting their size based on workload shifts \cite{Megiddo2003ARC}.

Building on these concepts, \textbf{CAR} (Clock with Adaptive
Replacement) adapts the self-tuning benefits of ARC to the more
efficient CLOCK structure \cite{Bansal2004CAR}. Similarly,
\textbf{CLOCK-Pro} introduces a ``test period'' for cold pages; items
are not promoted to the protected segment immediately upon insertion
but must demonstrate reuse within a test window
\cite{Jiang2005CLOCKPro}. While these algorithms improve scan
resistance, they rely on static rules that may not adapt quickly
enough to the complex, non-stationary patterns of modern AI and web
workloads \cite{Zhang2025Rethinking}.

\subsection{Machine Learning in Caching}
The complexity of modern access patterns has motivated the
integration of Machine Learning (ML) into cache management.
\textbf{LeCaR} (Learning Cache Replacement) employs Reinforcement
Learning (RL) using regret minimization to dynamically weight two
underlying policies (LRU and LFU) \cite{Vietri2018LeCaR}. While LeCaR
adapts to phase changes, it essentially selects between existing
static policies rather than predicting individual item utility.

Recent work has focused on ``lightweight'' ML to minimize inference
latency in real-time systems. Berger et al. demonstrated that
object-level features (e.g., size, type) could be used to predict hit
probabilities in CDNs \cite{Berger2018Towards}. More recently, the
importance of the \textit{promotion} path—deciding what enters the
cache—has gained attention. Chen et al. explored \textbf{lazy
promotion}, showing that deferring promotion decisions can
significantly reduce the churn caused by low-utility items
\cite{chen2025lazy}. Similarly, modern FIFO-based approaches like
\textbf{S3-FIFO} argue that simple queues with intelligent promotion
thresholds can outperform complex LRU variants \cite{Yang2023S3FIFO}.

\subsection{The PredCLOCK Approach}
Different from RL-based policy selection, \textbf{PredCLOCK} utilizes
a Supervised Learning model to act as an intelligent gatekeeper. By
predicting reuse probability upon access, we explicitly target the
inefficiency of blind promotion found in CLOCK and NbQ-CLOCK. This
approach aligns with recent findings that optimizing promotion
efficiency is often more impactful on hit ratios than optimizing
eviction policies \cite{chen2025lazy, Zhang2024SIEVE}.
